\documentclass[a4paper,11pt]{jsarticle}


% 数式
\usepackage{amsmath,amsfonts}
\usepackage{bm}
% 画像
\usepackage[dvipdfmx]{graphicx}


\begin{document}

\title{機械学習演習 14-問題 1}
\author{丸山 徹夜}
\date{\today}
\maketitle

\section*{1-1}

\begin{equation*}
  \frac{\partial{E}}{\partial{w}}
    = \frac{1}{2} \sum_{i=1}^{N} 2(w^T \phi(x_i) - y_i) \phi(x_i) + \frac{\alpha}{2}2w
    = \sum_{i=1}^{N} (w^T \phi(x_i) - y_i) \phi(x_i) + \alpha w
\end{equation*}

\begin{equation*}
  y = (y_1, y_2, \ldots, y_N)^Tとし、　
  \hat{y} = (w^T \phi(x_1), w^T \phi(x_2), \ldots, w^T \phi(x_N))^Tとすると、
\end{equation*}

\begin{equation*}
  \frac{\partial{E}}{\partial{w}} = \sum_{i=1}^{N} (\hat{y}_i - y_i) \phi(x_i) + \alpha w = \Phi^T (\hat{y} - y) + \alpha wに整理できる
\end{equation*}

\begin{equation*}
  \Phi^T (\hat{y} - y) + \alpha w = 0
\end{equation*}

\begin{equation*}
  \alpha wを右辺に移動し、左右両辺を -\alpha で割ると、
\end{equation*}

\begin{equation*}
  w = - \frac{1}{\alpha} \Phi^T (\hat{y} - y)
\end{equation*}

\begin{equation*}
  a = - \frac{1}{\alpha} (\hat{y} - y) とし、
\end{equation*}

\begin{equation*}
  w = \Phi^T a
\end{equation*}

\begin{equation*}
  \Rightarrow a_i = - \frac{1}{\alpha} (\hat{y}_i - y_i)
      = - \frac{1}{\alpha} (w^T \phi(x_i) - y_i)
\end{equation*}

\section*{1-2}
\subsection*{(a)}

\begin{equation*}
  \Phi = 
    \begin{pmatrix}
      \phi(x_1)^T \\
      \vdots \\
      \phi(x_N)^T 
    \end{pmatrix}、
  \Phi^T = 
    \begin{pmatrix}
      \phi(x_1) & \cdots & \phi(x_N)
    \end{pmatrix}
\end{equation*}

\begin{equation*}
  \Phi^T \Phi = 
    \begin{pmatrix}
      \phi(x_1) & \cdots & \phi(x_N)
    \end{pmatrix}
    \begin{pmatrix}
      \phi(x_1)^T \\
      \vdots \\
      \phi(x_N)^T 
    \end{pmatrix}
  = 
    \sum_{i=1}^{N} \phi(x_i) \phi(x_i)^T
\end{equation*}

\subsection*{(b)}
\begin{equation*}
  \sum_{i=1}^{N} a^T \Phi \phi(x_i) a^T \Phi \phi(x_i) = 
  \sum_{i=1}^{N} a^T \Phi \phi(x_i) (a^T \Phi \phi(x_i))^T = 
\end{equation*}

\begin{equation*}
  \sum_{i=1}^{N} a^T \Phi \phi(x_i) \phi(x_i)^T \Phi^T a =
  a^T \Phi (\sum_{i=1}^{N} \phi(x_i) \phi(x_i)^T) \Phi^T a =
  a^T \Phi \Phi^T \Phi \Phi^T a
\end{equation*}

\subsection*{(c)}

\begin{equation*}
  \sum_{i=1}^{N} \phi(x_i)y_i =
    (\phi(x_1), \ldots, \phi(x_N))
    \begin{pmatrix}
      y_1 \\
      \vdots \\
      y_N
    \end{pmatrix}
  = \Phi^T y
\end{equation*}

\subsection*{(d)}

\begin{equation*}
  \sum_{i=1}^{N} y_i^2 =
    (y_1, \ldots, y_N)
    \begin{pmatrix}
      y_1 \\
      \vdots \\
      y_N
    \end{pmatrix}
  = y^T y
\end{equation*}

\section*{1-3}
\begin{equation*}
  E(\Phi^T a)
    = \frac{1}{2} \sum_{i=1}^{N} ((\Phi^T a)^T \phi(x_i) - y_i)^2 + \frac{\alpha}{2} (\Phi^T a)^T \Phi^T a
    = \frac{1}{2} \sum_{i=1}^{N} (a^T \Phi \phi(x_i) - y_i)^2 + \frac{\alpha}{2} a^T \Phi \Phi^T a
\end{equation*}
\begin{equation*}
  = \frac{1}{2} \sum_{i=1}^{N} (a^T \Phi \phi(x_i) a^T \Phi \phi(x_i) - 2a^T \Phi \phi(x_i) y_i + y_i^2) + \frac{\alpha}{2} a^T \Phi \Phi^T a
\end{equation*}
\begin{equation*}
  E(a) = \frac{1}{2} a^T \Phi \Phi^T \Phi \Phi^T a - a^T \Phi \Phi^T y + \frac{1}{2} y^T y + \frac{\alpha}{2} a^T \Phi \Phi^T a
\end{equation*}

\section*{1-4}
\begin{equation*}
  K = \Phi \Phi^Tとして、
\end{equation*}
\begin{equation*}
  E(a) = \frac{1}{2} a^T K K a - a^T K y + \frac{1}{2} y^T y + \frac{\alpha}{2} a^T K a
\end{equation*}

\section*{1-5}
\begin{equation*}
  a_i = - \frac{1}{\alpha} (w^T \phi(x_i) - y_i) \Rightarrow a = - \frac{1}{\alpha} (\Phi w - y)
\end{equation*}

$wに \Phi^T a を代入すると、$
\begin{equation*}
  a = - \frac{1}{\alpha} (\Phi \Phi^T a - y)
\end{equation*}
$aについて整理すると、$
\begin{equation*}
  (K + \alpha I_N) a = y
\end{equation*}
$左右両辺に(K + \alpha I_N)^{-1}をかけると$
\begin{equation*}
  a = (K + \alpha I_N)^{-1} y
\end{equation*}

\section*{1-6}

双対問題のパラメータベクトルλに対するラグランジュ関数がハードマージンSVMの場合、ソフトマージン SVMの場合ともに下記の関数で表します。
\begin{equation*}
  L(\lambda) = - \frac{1}{2} \sum_{i \in SV} \sum_{j \in SV} \lambda_i \lambda_j t_i t_j x_i^T x_j + \sum_{i \in SV} \lambda_i
\end{equation*}

$i, j \notin SVの場合、λ_i = 0, λ_j = 0 $

\begin{equation*}
  L(\lambda) = - \frac{1}{2} \sum_{i=1}^{N} \sum_{j=1}^{N} \lambda_i \lambda_j t_i t_j x_i^T x_j + \sum_{i=1}^{N} \lambda_i
\end{equation*}

$x_i, x_jを、\phi(x_i), \phi(x_j)にそれぞれ置き換えると、$

\begin{equation*}
  L(\lambda)
    = - \frac{1}{2} \sum_{i=1}^{N} \sum_{j=1}^{N} \lambda_i \lambda_j t_i t_j \phi(x_i)^T \phi(x_j) + \sum_{i=1}^{N} \lambda_i
    = - \frac{1}{2} \sum_{i=1}^{N} \sum_{j=1}^{N} \lambda_i \lambda_j t_i t_j k(x_i, x_j) + \sum_{i=1}^{N} \lambda_i
\end{equation*}

線形識別関数：
\begin{equation*}
  f(x)
    = \sum_{i \in SV} \lambda_i t_i x_i^T x + b
    = \sum_{i=1}^{N} \lambda_i t_i \phi(x_i)^T \phi(x) + b
    = \sum_{i=1}^{N} \lambda_i t_i k(x_i, x) + b
\end{equation*}

\end{document}

